% !TEX root = ../main.tex
\chapter{Federal Funds, Final Settlement}\label{ch:fedfunds}
\begin{paracol}{2}

\begin{sources}
\begin{tabularx}{\linewidth}{@{}l l L@{}}
\toprule
\textbf{Segment} & \textbf{Length} & \textbf{Content}\\
\midrule
\seg{6}{1}{L6.1 FT: European Bank Deleveraging} & 5:43 & European banks sell property loans to meet capital ratios; three overnight rates in the FT.\\
\seg{6}{2}{L6.2 What are Fed Funds?} & 5:13 & An interbank market in reserves; the Fed funds rate over the cycle; excess reserves.\\
\seg{6}{3}{L6.3 Payment settlement versus Required Reserves} & 1:11 & The two jobs of a Fed funds desk (Stigum).\\
\seg{6}{4}{L6.4 Payment elasticity/discipline, public and private} & 9:03 & Daylight overdrafts at the Fed; CHIPS; real-time gross settlement.\\
\seg{6}{5}{L6.5 The Function of the Fed Funds Market} & 9:15 & Bilateral overnight borrowing replaces central overnight lending; paying today by promising tomorrow.\\
\seg{6}{6}{L6.6 Payment versus Funding: an example} & 11:24 & Buying an apartment with a mortgage: the Fed funds market and a matched-book dealer.\\
\seg{6}{7}{L6.7 Brokers versus Dealers} & 2:03 & The difference.\\
\seg{6}{8}{L6.8 Payments Imbalances and the Fed Funds Rate} & 7:06 & Imbalances, dealer prices, and how the Fed influences the rate through repo.\\
\seg{6}{9}{L6.9 Secured versus Unsecured Interbank Credit} & 5:05 & Fed funds are unsecured; credit lines; stress in unsecured markets.\\
\seg{6}{10}{L6.10 Required Reserves, redux} & 7:20 & Reserve requirements, lagged and contemporaneous reserve accounting, and why they are ineffective.\\
\bottomrule
\end{tabularx}

\smallskip
\textbf{Files.} Transcripts \file{materials/transcripts/L06-P*.txt}.
\textbf{Reading.} Stigum, \emph{The Money Market}, the chapter on the Fed funds market (quoted in class).
\end{sources}

\section*{Motivation}
The next three lectures cover the three main dollar money markets \ts{6}{1}{3:57}\ts{6}{1}{4:18}:
\begin{enumerate}
  \item \textbf{Fed funds} (this lecture): unsecured interbank loans of reserves at the Fed, the market for ``the best money in the world'';
  \item \textbf{repo} (lecture~7): secured loans;
  \item \textbf{Eurodollars} (lecture~8): unsecured interbank loans of dollars offshore.
\end{enumerate}
In the FT of the day their overnight rates are 0.15\% (US dollar LIBOR, the Eurodollar rate), 0.29\% (US overnight repo) and 0.16\% (Fed funds
effective): all overnight, all different. Why are they what they are? \ts{6}{1}{4:38}\ts{6}{1}{5:00}\ts{6}{1}{5:23} The thread of this lecture: the
Fed funds market exists to meet the \emph{survival constraint} of banks at the end of each day, i.e.\ \emph{final settlement}.

% ------------------------------------------------------------------
\section{FT: European bank deleveraging}\label{sec:ff-ft}

\begin{casestudy}[FT, September 2012: ``Europe's banks race to offload property loans'']
European banks hold property loans left over from the boom in Spain, Ireland and elsewhere, many no longer worth face value
\ts{6}{1}{0:28}\ts{6}{1}{0:50}. A new round of regulation will require more capital against such risky loans, so that losses fall on private
capital rather than the taxpayer \ts{6}{1}{1:10}\ts{6}{1}{1:31}. Capital requirements are ratios of capital to risky assets; rather than raise
capital, the banks shrink the denominator: they package the loans, securitise some, and sell them, sometimes ``at pennies on the dollar'', to
private equity funds, and use the cash to pay down short-term borrowing or borrowing from the ECB \ts{6}{1}{1:52}\ts{6}{1}{2:14}\ts{6}{1}{2:56}.
\end{casestudy}
\begin{taccts}
\tacct[2.4cm]{European bank}{$-$Property loans\\$+$Cash, then\\$-$Cash}{then: $-$Short-term borrowing}\qquad
\tacct[2.4cm]{Private equity fund}{$+$Property loans\\$-$Cash}{}
\end{taccts}
This is \emph{deleveraging}: shrinking the balance sheet on both sides to reach the desired capital ratio \ts{6}{1}{2:56}. The loans are sold at
fire-sale prices, and the banks make no new loans: a credit crunch in Europe. The article is about the connection between banks and other
financial institutions, and between banking and the real economy \ts{6}{1}{3:16}\ts{6}{1}{3:37}.

\begin{definition}[New concepts: capital requirement, securitisation]
A \emph{capital requirement} sets a minimum ratio of a bank's capital (net worth) to its (risk-weighted) assets. \emph{Securitisation} is the pooling
of loans into a package whose cash flows back securities that can be sold to investors (the loans are turned into ``securities'').
\end{definition}

% ------------------------------------------------------------------
\section{What are Fed funds?}\label{sec:ff-what}

In the flow of funds matrix for 2011 (here in stock form, assets and liabilities), line 12, ``Fed funds and security repos'', shows for the
domestic financial sectors \$823 billion of assets and \$1.1 trillion of liabilities: an \emph{interbank} market, in which the financial sector
borrows from and lends to itself; the numbers differ because some non-banks participate too \ts{6}{2}{0:07}\ts{6}{2}{0:28}\ts{6}{2}{0:52}\ts{6}{2}{1:12}.

\begin{definition}[New concepts: Fed funds]\label{def:ff}
\emph{Federal funds} are unsecured loans, typically overnight, of reserves (deposits at the Federal Reserve) between banks: a promise to pay reserves
tomorrow morning in exchange for reserves today \ts{6}{5}{5:17}. Fed funds are \emph{not} a liability or an asset of the Fed, despite the name: they are
assets and liabilities of banks \ts{6}{2}{2:34}.
\end{definition}

In the hierarchy, from the banks' point of view reserves are money and Fed funds are credit: interbank promises to pay money, a credit expansion on the
base of the money on the Fed's balance sheet \ts{6}{2}{1:53}\ts{6}{2}{2:13}. ``If you get off on the wrong foot there, you can be halfway through the lecture
before realizing where you are'' \ts{6}{2}{2:34}.

\paragraph{The rate.} The Fed funds effective rate (FRED, St.~Louis Fed) fluctuates with the business cycle; in the crisis the Fed lowered its target from 5\%
to 2\% and then to nearly zero, where it has stayed, in a range of 0 to 0.25\%, with the effective rate at 0.16\% \ts{6}{2}{2:54}\ts{6}{2}{3:16}\ts{6}{2}{3:36}. The
Fed used to steer the economy by moving this rate; at the zero bound it uses instead the \emph{size} of its balance sheet, quantitative easing
\ts{6}{2}{3:57}\ts{6}{2}{4:17}. Reserves used to be about \$50 billion; now there are over a trillion, excess reserves: reserves are not scarce any more, and yet the
Fed funds market is still active, because it serves the payment system \ts{6}{2}{4:38}\ts{6}{2}{4:58}.

\begin{finance}[The two jobs of a Fed funds desk]
Stigum: the primary job of the manager of a bank's Fed funds desk is to ensure
\begin{enumerate*}[(i)]
  \item that the bank settles with the Fed, every day, and
  \item that in doing so it holds no more excess reserves than it can carry into the next week
\end{enumerate*} \ts{6}{3}{0:09}. The second job, tied to reserve requirements, is now basically defunct: with so many reserves, missing a reserve target is
``ridiculous''. Only the first, the payment job, matters \ts{6}{3}{0:30}\ts{6}{3}{0:50}.
\end{finance}

% ------------------------------------------------------------------
\section{Intraday elasticity: daylight overdrafts and CHIPS}\label{sec:ff-intraday}

The survival constraint, in Stigum's words that the bank must settle with the Fed, takes the precise institutional form of a requirement to \emph{end
the day with a non-negative balance at the Fed} \ts{6}{4}{0:08}. It does not bind \emph{during} the day \ts{6}{4}{0:29}.

\paragraph{Daylight overdrafts at the Fed.} Mehrling spent the summer of 1986 at the Board of Governors as a graduate student, when the Fed was worried
about its credit exposure from daylight overdrafts \ts{6}{4}{0:49}\ts{6}{4}{1:10}. Bank A must pay bank B but has no reserves: it simply runs an overdraft,
which is a loan from the Fed; the Fed creates the reserves and credits them to B, which feels paid and knows nothing of the arrangement
\ts{6}{4}{1:33}\ts{6}{4}{1:53}\ts{6}{4}{2:14}\ts{6}{4}{2:34}:
\begin{taccts}
\tacct[2.2cm]{Fed}{$+$Overdraft of A}{$+$Reserves of B}\quad
\tacct[2.2cm]{Bank A}{}{$-$Deposits\\$+$Overdraft at Fed}\quad
\tacct[2.2cm]{Bank B}{$+$Reserves}{$+$Deposits}
\end{taccts}
This is ``the credit grease that makes the payment system go during the day'' \ts{6}{4}{2:54}. The overdraft is a promise to come up with the reserves by
the end of the day; if A does not, the Fed is forced to lend overnight, but at a penalty, 100 basis points over the Fed funds rate as said in class, to give
A an incentive to find someone who will lend cheaper. The Fed cannot say ``pay me or you're dead'', because it is protecting the payment system; it says
``pay me or you owe me more than you want to''. Sometimes a bank cannot find a counterparty and pays the penalty \ts{6}{4}{3:15}\ts{6}{4}{3:36}\ts{6}{4}{3:56}\ts{6}{4}{4:16}.
Credit expands during the day and collapses at the end of it: the intraday version of the pattern at every level and every time scale \ts{6}{4}{4:36}\ts{6}{4}{4:57}.

\paragraph{CHIPS.} The same happens on the balance sheet of a private clearinghouse, CHIPS, a mutual organisation of banks that collects dues to and from
during the day and settles at the end of the day on Fedwire \ts{6}{4}{4:57}\ts{6}{4}{5:23}\ts{6}{4}{5:44}. Every morning members deposit collateral; during the day
A can owe CHIPS while B is owed by CHIPS \ts{6}{4}{6:05}\ts{6}{4}{6:25}\ts{6}{4}{6:46}\ts{6}{4}{7:07}:
\begin{taccts}
\tacct[2.2cm]{CHIPS}{Due from A}{Due to B}\quad
\tacct[2.2cm]{Bank A}{}{Due to CHIPS}\quad
\tacct[2.2cm]{Bank B}{Due from CHIPS}{}
\end{taccts}
Two layers of intraday credit expansion, a public one on the Fed's balance sheet and a private one on CHIPS's \ts{6}{4}{6:05}. If A fails, all members of CHIPS
are supposed to pay B: B holds an obligation of the clearinghouse, better than one of A, though not as good as one of the Fed \ts{6}{4}{7:28}\ts{6}{4}{7:49}.

\begin{definition}[New concepts: real-time gross settlement]
A \emph{real-time gross settlement} (RTGS) system settles each payment individually (gross) and immediately (real time) in central bank money: once
received, the reserves are money \emph{now}, not to be netted at the end of the day. Fedwire is RTGS; CHIPS nets during the day and settles at the end,
so a CHIPS credit is money only ``hopefully'', after settlement on Fedwire \ts{6}{4}{8:09}\ts{6}{4}{8:30}\ts{6}{4}{8:50}.
\end{definition}

\begin{coursenote}[Daylight overdraft policy]
The Fed and CHIPS have tightened up over time: the size of overdrafts allowed depends on a bank's capital, and the Fed charges for them, on the average
overdraft during the day, so a trillion-dollar overdraft for a few microseconds costs essentially nothing, but still disciplines the bank, which must know
where the money is coming from \ts{6}{5}{2:52}\ts{6}{5}{3:13}\ts{6}{5}{3:33}. (Fees on daylight overdrafts were introduced in 1994, with caps related to capital.
The overnight penalty rate quoted in class may have changed since.) Large security and international transactions cluster at the end of the day, when CHIPS
settles through Fedwire \ts{6}{5}{3:53}.
\end{coursenote}

% ------------------------------------------------------------------
\section{The function of the Fed funds market}\label{sec:ff-function}

At the end of the day, in a closed system, the overdrafts and surpluses sum to zero: if I owe the Fed, somebody out there has excess reserves, and if we
find each other we can both end flat and neither needs the Fed; the same for CHIPS \ts{6}{5}{0:48}\ts{6}{5}{1:09}\ts{6}{5}{1:30}. \emph{Interbank markets exist to let
banks find each other} \ts{6}{5}{1:51}. For CHIPS the relevant interbank market is the Eurodollar market: anyone can be a member of CHIPS, or have a
correspondent that is, including many foreign banks; Fedwire is a much more exclusive club, and its members borrow and lend in the Fed funds market.
Roughly, Fed funds are domestic, public, ``Fed'' dollars; Eurodollars international, private dollars; both serve the same function
\ts{6}{5}{2:11}\ts{6}{5}{2:32}. The Fed funds market creates bilateral overnight borrowing to substitute for the centralised overnight lending the Fed would
otherwise be forced to do \ts{6}{5}{4:14}.

\paragraph{A Fed funds trade.} Bank A, short at the end of the day, borrows from bank B, long: it promises reserves tomorrow morning and gets reserves today,
which pay off its overdraft; the Fed's balance sheet expands during the day and contracts back \ts{6}{5}{4:35}\ts{6}{5}{4:57}\ts{6}{5}{5:37}\ts{6}{5}{5:58}\ts{6}{5}{6:19}.
\begin{taccts}
\tacct[2.4cm]{Bank A}{$+$Reserves (repays overdraft)}{$+$Fed funds borrowed}\qquad
\tacct[2.4cm]{Bank B}{$-$Reserves\\$+$Fed funds lent}{}
\end{taccts}
Banks do not wait for the end of the day: a desk anticipates needs and trades all day long \ts{6}{5}{6:40}\ts{6}{5}{7:00}. Without any overdraft: A borrows reserves
from B and then uses them to pay B; the reserves go back and forth and what is left is a Fed funds loan \ts{6}{5}{7:21}\ts{6}{5}{8:02}\ts{6}{5}{8:30}:
\begin{taccts}
\tacct[2.4cm]{Bank A}{}{$-$Deposits of payer\\$+$Fed funds borrowed}\qquad
\tacct[2.4cm]{Bank B}{$+$Fed funds lent}{$+$Deposits of payee}
\end{taccts}

\begin{proposition}[Paying today by promising tomorrow]
A Fed funds loan pushes final settlement into the future: ``I'm going to pay you today by promising to pay you tomorrow'' \ts{6}{5}{8:50}\ts{6}{5}{9:11}.
\end{proposition}

% ------------------------------------------------------------------
\section{Payment versus funding: an example}\label{sec:ff-example}

Mehrling buys your apartment in Manhattan for \$1 million, with a mortgage from his bank, Citibank; you bank at Chase. Credit checks are ``interesting but not
monetary economics'' \ts{6}{6}{0:08}\ts{6}{6}{0:28}\ts{6}{6}{0:49}. Step by step:
\begin{enumerate}
  \item \textbf{The loan.} Citibank grants the mortgage by crediting a deposit: a swap of IOUs, but with different timing and different signers: Mehrling
        owes \$1 million over 30 years, Citibank owes \$1 million on demand; when Citibank signs an IOU it is money, when Mehrling does it is credit
        \ts{6}{6}{1:09}\ts{6}{6}{1:29}\ts{6}{6}{1:49}\ts{6}{6}{2:10}.
  \item \textbf{Funding the payment.} The deposit must go to Chase, and Citibank does not keep \$1 million of idle reserves. Its credit department knows the
        closing is coming, so it borrows \$1 million in the Fed funds market, not from Chase but from a third party, HSBC, which lends reserves it has (or
        runs an overdraft) \ts{6}{6}{2:30}\ts{6}{6}{2:51}\ts{6}{6}{3:34}\ts{6}{6}{3:54}. Eventually the mortgage will be packaged and securitised, but that takes weeks; the
        payment must happen now \ts{6}{6}{3:11}.
  \item \textbf{The payment.} Citibank pays Chase in reserves and Chase credits your deposit: it is happy to, because it is settled, it has money
        \ts{6}{6}{4:14}\ts{6}{6}{4:35}\ts{6}{6}{4:57}\ts{6}{6}{5:17}.
  \item \textbf{Chase lends on.} Reserves pay nothing, so Chase lends them in the Fed funds market, to HSBC \ts{6}{6}{5:38}\ts{6}{6}{5:59}.
\end{enumerate}
\begin{taccts}
\tacct[2.0cm]{Citibank}{$+$Mortgage}{$+$Fed funds (HSBC)}\quad
\tacct[2.0cm]{HSBC}{$+$Fed funds (Citi)}{$+$Fed funds (Chase)}\quad
\tacct[2.0cm]{Chase}{$+$Fed funds (HSBC)}{$+$Deposit (you)}

\medskip
\tacct[2.2cm]{Mehrling}{$+$Apartment}{$+$Mortgage}\qquad
\tacct[2.2cm]{You}{$-$Apartment\\$+$Deposit}{}
\end{taccts}

\paragraph{Who lent the money?} Follow the chain: Citibank borrowed from HSBC, HSBC from Chase, Chase from you. ``In a certain sense I bought it from you and you
lent me the money for it'', though Mehrling owes the intermediaries, not you \ts{6}{6}{6:19}\ts{6}{6}{6:40}\ts{6}{6}{7:00}. Until the mortgage is financed by some
pension fund, it has been funded through the \emph{payment system}, by an expansion of balance sheets: nobody saved extra money so that Mehrling could borrow;
you transformed a house into a deposit, he acquired a house and a mortgage \ts{6}{6}{7:00}\ts{6}{6}{7:21}.

\paragraph{HSBC as a dealer.} HSBC lends and borrows Fed funds at the same time, flat in exposure: it does so only because it lends at a higher rate than it
borrows. It is a \emph{dealer} in Fed funds, quoting rates at which it will lend and borrow and setting them so that over the day inflows match outflows; when
they do not, it must find funds itself, or place the surplus \ts{6}{6}{7:41}\ts{6}{6}{8:01}\ts{6}{6}{8:22}\ts{6}{6}{8:44}.

\begin{definition}[New concepts: matched book]
A dealer runs a \emph{matched book} when its borrowing and its lending in a market have the same size (and maturity), so that it earns the spread without
net exposure to the price \ts{6}{6}{8:44}\ts{6}{6}{9:04}.
\end{definition}

\paragraph{An internal drain.} If you withdraw the \$1 million in cash (``1MM'', in Stigum's notation), Chase must get currency: it asks the Fed to convert
reserves into cash, which is physically ``downtown at 33 Liberty Street''. The Fed's liabilities change form, one million more of currency and one million less
of reserves; an internal drain, which causes no monetary problem, since both are liabilities of the Fed \ts{6}{6}{9:25}\ts{6}{6}{10:06}\ts{6}{6}{10:46}\ts{6}{6}{11:07}.

% ------------------------------------------------------------------
\section{Brokers versus dealers}\label{sec:ff-brokers}

A \emph{broker} puts buyer and seller together and takes a fee, like a real-estate broker, who does not buy your house and resell it. In a brokered Fed funds
trade the loan sits on the balance sheets of A and B only \ts{6}{7}{0:07}\ts{6}{7}{0:28}\ts{6}{7}{0:48}. A \emph{dealer} like HSBC stands in between, and the loans sit on
its balance sheet: Chase and Citibank probably do not even know that they are, in effect, lending to and borrowing from each other \ts{6}{7}{1:08}\ts{6}{7}{1:29}. The
dealer is ``the source of liquidity in these markets'', a theme for the rest of the course \ts{6}{7}{1:50}. (Brokers and dealers were defined in
\cref{def:nh-mm}.)

% ------------------------------------------------------------------
\section{Payment imbalances and the Fed funds rate}\label{sec:ff-rate}

During the day the Fed and CHIPS expand to make payments possible; at the end both shrink back, and the Fed funds market is one way to do it. But the credit does
not vanish: it moves to the bilateral balance sheets of borrowers and lenders and onto the balance sheets of the Fed funds dealers \ts{6}{8}{0:07}\ts{6}{8}{0:28}\ts{6}{8}{0:50}.
The centre protects itself (``we won't extend this overnight, you do it''); daily payment imbalances that could not be cleared show up as expansion of Fed funds credit
and of dealer balance sheets, pushed into the future in the hope that tomorrow the flow reverses \ts{6}{8}{1:11}\ts{6}{8}{1:32}\ts{6}{8}{1:52}. These quantities are highly
sensitive measures of payment patterns, and so are the prices: a dealer with many more borrowers than lenders raises its rate to attract lenders, and lowers it in the
opposite case. The Fed funds rate is a \emph{market rate}, and it fluctuates within the day too \ts{6}{8}{1:52}\ts{6}{8}{2:14}\ts{6}{8}{2:34}\ts{6}{8}{2:56}.

\paragraph{How does the Fed ``set'' the rate?} Not as in intermediate macro: the Fed does not participate in the Fed funds market at all, neither borrowing nor
lending \ts{6}{8}{3:16}\ts{6}{8}{3:37}. It has a \emph{target}, and it influences the rate by changing the quantity of reserves, by entering other markets, in particular the
repo market \ts{6}{8}{3:57}. Almost every day it anticipates payment imbalances that would push the rate away from target, and adjusts the quantity of reserves: it
changes money in order to influence credit \ts{6}{8}{4:18}\ts{6}{8}{4:38}\ts{6}{8}{4:59}. To add reserves it makes an overnight repurchase agreement loan, for example to a
security dealer funding its securities, which ends up with reserves it does not want and lends them in the Fed funds market, putting downward pressure on the rate
\ts{6}{8}{5:20}\ts{6}{8}{5:41}\ts{6}{8}{6:01}\ts{6}{8}{6:21}:
\begin{taccts}
\tacct[2.0cm]{Fed}{$+$Repo loan}{$+$Reserves}\quad
\tacct[2.0cm]{Security dealer}{$+$Reserves, then $-$Reserves\\$+$Fed funds lent}{$+$Repo from Fed}
\end{taccts}
The \emph{Fed funds effective} rate is the average of the rates at which funds trade during the day; good borrowers pay less, bad ones more \ts{6}{8}{6:45}.

\begin{definition}[New concepts: target rate, effective rate]
The \emph{Fed funds target} is the rate (since 2008 a range) that the Federal Open Market Committee (FOMC) aims at; the \emph{effective} Fed funds rate is the
volume-weighted average of the rates actually paid in the market on a given day.
\end{definition}

% ------------------------------------------------------------------
\section{Secured versus unsecured interbank credit}\label{sec:ff-unsecured}

Fed funds are \emph{unsecured} overnight loans: no collateral. If Citibank does not pay HSBC, HSBC can sue, but there is no asset to seize; a mortgage, by contrast,
is secured by the house \ts{6}{9}{0:08}\ts{6}{9}{0:29}\ts{6}{9}{0:49}\ts{6}{9}{1:10}. Eurodollars too are unsecured interbank loans \ts{6}{9}{1:10}. Why lend a million dollars
overnight on a promise? Because banks are in a network, it is their business, and tomorrow they will be on the other side; they watch their exposure to each
counterparty very closely, with \emph{credit lines}: ``we're already in them for \$10 million, tell them your credit line to us is full'', and they diversify
their lending \ts{6}{9}{1:30}\ts{6}{9}{1:50}\ts{6}{9}{2:11}\ts{6}{9}{2:32}.

The market, especially the Eurodollar market, is under great stress: with many insolvent banks in Europe, nobody wants to lend unsecured, and even a dealer in the
middle worries, because its business is highly leveraged on both sides: if Citibank does not pay HSBC, HSBC must still pay Chase; it is not a pass-through, it faces its
own survival constraint \ts{6}{9}{2:53}\ts{6}{9}{3:13}\ts{6}{9}{3:33}\ts{6}{9}{3:54}. Repurchase agreements, next lecture, are \emph{secured}, by a Treasury bill for example
\ts{6}{9}{3:54}. This is where large amounts are transferred, where risk is taken and liquidity created: the Fed funds market is the market for the best money in
the world, Eurodollars another unsecured market, repo the secured one; understanding them puts us at the very core of the international payment system
\ts{6}{9}{4:14}\ts{6}{9}{4:35}\ts{6}{9}{4:56}.

\begin{definition}[New concepts: secured and unsecured credit, credit line]
A \emph{secured} loan is backed by collateral that the lender can take if the borrower defaults; an \emph{unsecured} loan rests only on the borrower's general
promise. A \emph{credit line} is the maximum exposure a lender is willing to have to a given counterparty.
\end{definition}

% ------------------------------------------------------------------
\section{Required reserves, redux}\label{sec:ff-required}

``The least important part of this lecture'': reserve requirements \ts{6}{10}{0:07}. Stigum describes two-week \emph{maintenance periods} over which a bank must hold
an average quantity of reserves, so a bank had two tasks: settle every day (easy, with the Fed funds market) and hit the average without overshooting
\ts{6}{10}{0:27}\ts{6}{10}{0:48}.

\paragraph{Monetarism and reserve accounting.} Under Volcker, monetarists wanted to use reserve requirements to control the money supply: if deposits require
reserves, a shortage of reserves limits the expansion of deposits and of bank credit \ts{6}{10}{1:08}\ts{6}{10}{1:29}.
\begin{itemize}
  \item With \emph{lagged reserve accounting}, required reserves depended on deposits of two weeks earlier: the loan and the deposit are already there, so the
        Fed has to supply the reserves, or somebody will miss the requirement; the Fed only accommodates demand \ts{6}{10}{1:50}\ts{6}{10}{2:12}\ts{6}{10}{2:32}.
  \item With \emph{contemporaneous reserve accounting}, required reserves depended on deposits in the same period: harder to hit, with both deposits and
        reserves moving; it caused a lot of chaos in the Fed funds market and was abandoned \ts{6}{10}{2:53}\ts{6}{10}{3:15}\ts{6}{10}{3:36}.
\end{itemize}
It was an attempt to impose discipline on a system that seemed out of control, in a period of double-digit inflation; according to Stigum it did not really work
\ts{6}{10}{3:36}\ts{6}{10}{3:56}. Today, with trillions of reserves, requirements are a dead letter, yet the Fed funds market exists for payments \ts{6}{10}{4:17}.

\begin{coursenote}[Dates]
The Fed's ``monetarist experiment'' of targeting reserves ran from October 1979 to 1982. Contemporaneous reserve accounting was introduced in February 1984 and
replaced by lagged accounting again in 1998. (In March 2020, after these lectures, the Fed set reserve requirement ratios to zero.)
\end{coursenote}

\paragraph{Were reserve requirements ever important?} Asked whether they will matter again after the crisis, Mehrling answers that a trillion of excess reserves will
not go away soon, but the deeper question is whether the quantity of reserves ever significantly restricted credit or money \ts{6}{10}{4:38}\ts{6}{10}{4:59}\ts{6}{10}{5:19}.
A parallel banking system, shadow banking, creates close substitutes for money and credit with no reserve or capital requirement, possibly offshore; marginal credit can flow
there. Squeezing banks, as with the European capital requirements of the FT article, does not mean there is no credit, only that it does not come from banks: ``like a
balloon, you punch it here and it blows out there'' \ts{6}{10}{5:39}\ts{6}{10}{5:59}\ts{6}{10}{6:20}. Mehrling sides with those who think the reserve constraint is ineffective in
modern financial markets; it probably was not in 1950, when it was invented and the textbooks were written, because there were no alternatives \ts{6}{10}{6:40}\ts{6}{10}{7:00}.

% ------------------------------------------------------------------
\flushside
\section*{Key formulas and ideas}
\begin{keyformulas}
\begin{tabularx}{\linewidth}{@{}l L@{}}
Fed funds & unsecured overnight interbank loans of reserves; assets/liabilities of banks, not of the Fed\\
Survival constraint & end the day with a non-negative balance at the Fed\\
Intraday & daylight overdrafts at the Fed (RTGS) and at CHIPS (net): credit expansion during the day\\
End of day & surplus and deficit banks find each other in the Fed funds (or Eurodollar) market; else penalty borrowing from the Fed\\
Payment vs funding & a payment funded by interbank borrowing: paying today by promising tomorrow\\
Dealer & borrows and lends at a spread, matched book; source of market liquidity\\
Fed funds rate & a market rate; the Fed targets it by adjusting reserves (repo), not by trading Fed funds\\
Reserve requirements & lagged vs contemporaneous accounting; ineffective with shadow banking\\
\end{tabularx}
\end{keyformulas}

\section*{Exercises}

\begin{exercise}[A day at the Fed]
Bank A starts the day with 0 reserves, B with 50. During the day A pays B 80 and B pays A 30.
\begin{enumerate}[(a)]
  \item Track A's balance at the Fed; what is its maximum daylight overdraft?
  \item What must A do by the end of the day? Write the T-accounts if it borrows from B in the Fed funds market.
  \item Repeat when B is unwilling to lend to A, and a dealer D borrows from B and lends to A.
\end{enumerate}
\end{exercise}

\begin{exercise}[The apartment]
Redo the apartment example when Chase itself lends Citibank the reserves in the Fed funds market. Then suppose the buyer's bank and the seller's bank are the same:
what happens to reserves and to the Fed funds market?
\end{exercise}

\begin{exercise}[Dealer pricing]
A Fed funds dealer quotes 0.15\% to borrow and 0.20\% to lend. During the morning it receives requests to borrow 300 and offers to lend 100.
\begin{enumerate}[(a)]
  \item What is its net position if it accepts everything?
  \item Which way should it move its quotes, and why does this move the effective rate?
\end{enumerate}
\end{exercise}

\begin{exercise}[Reserve accounting]
A bank with deposits of 1000 and a 10\% reserve requirement grants a loan of 100 that is spent within the bank. Compare what it must do under lagged and under
contemporaneous reserve accounting, and explain why the Fed must accommodate in the first case.
\end{exercise}
